% =============================================================================
% preambulo/pacotes.tex  --  Pacotes adicionais sobre a classe USPSC
%
%   A classe USPSC já carrega: memoir/abntex2 base, geometry, fancyhdr,
%   xcolor (com 'capa-azul'), hyperref/url, microtype, textpos, ...
%   Aqui adicionamos apenas o que é específico do TCC.
% =============================================================================

% --- Codificação e fonte ---
\usepackage[T1]{fontenc}
\usepackage[utf8]{inputenc}
\usepackage{lmodern}            % Fonte Latin Modern (padrão recomendado USPSC)
\usepackage{lastpage}           % Usado pela ficha catalográfica
\usepackage{indentfirst}        % Indenta primeiro parágrafo de cada seção
\usepackage{color}              % Controle de cores
\usepackage{graphicx}           % Inclusão de figuras
\usepackage{float}              % Posicionamento estrito de figuras/tabelas

% --- Pacotes requeridos pela capa ICMC do Pacote USPSC ---
\usepackage{microtype}          % Melhorias de justificação
\usepackage{pdfpages}           % Para \includepdf (folha de aprovação, ficha)
\usepackage{makeidx}            % Índice remissivo
\usepackage{hyphenat}           % Comando \nohyphens da capa ICMC
\usepackage[absolute]{textpos}  % Comando \begin{textblock*} da capa ICMC
\usepackage{eso-pic}            % \AddToShipoutPicture / background da capa
\usepackage{makebox}            % Caixas de texto na capa

% --- Citações ABNT (Sistema autor-data) ---
% Empregamos o estilo alfabético do Pacote USPSC, compatibilizado com
% ABNT NBR 6023:2018 e 10520:2023 via abntex2-alf-USPSC.bst.
\usepackage[
  alf,
  abnt-emphasize=bf,
  abnt-thesis-year=both,
  abnt-repeated-author-omit=no,
  abnt-last-names=abnt,
  abnt-etal-cite=3,
  abnt-etal-list=3,
  abnt-etal-text=it,
  abnt-and-type=e,
  abnt-doi=doi,
  abnt-url-package=none,
  abnt-verbatim-entry=no
]{abntex2cite}
\bibliographystyle{USPSC-classe/abntex2-alf-USPSC}

% --- Nota de rodapé ---
\renewcommand{\footnotesize}{\small}

% --- Tabelas ---
\usepackage{multicol}
\usepackage{multirow}
\usepackage{longtable}
\usepackage{threeparttablex}
\usepackage{array}
\usepackage{booktabs}
\usepackage{pdflscape}             % páginas em paisagem (Apêndice A)           % Tabelas profissionais (toprule/midrule)

% --- Listagens de código (será útil em apêndices) ---
\usepackage{listings}
\usepackage{xcolor}             % já carregado pela USPSC; reincluir não dói
\lstset{
  basicstyle=\ttfamily\small,
  breaklines=true,
  columns=fullflexible,
  showstringspaces=false,
  frame=single,
  numbers=left,
  numberstyle=\tiny\color{gray},
  keywordstyle=\bfseries,
  commentstyle=\itshape\color{gray!70!black},
  tabsize=2,
  extendedchars=true,
  inputencoding=utf8,
  literate=%
    {á}{{\'a}}1 {é}{{\'e}}1 {í}{{\'i}}1 {ó}{{\'o}}1 {ú}{{\'u}}1%
    {Á}{{\'A}}1 {É}{{\'E}}1 {Í}{{\'I}}1 {Ó}{{\'O}}1 {Ú}{{\'U}}1%
    {à}{{\`a}}1 {è}{{\`e}}1 {ì}{{\`i}}1 {ò}{{\`o}}1 {ù}{{\`u}}1%
    {ã}{{\~a}}1 {õ}{{\~o}}1 {ñ}{{\~n}}1%
    {Ã}{{\~A}}1 {Õ}{{\~O}}1%
    {â}{{\^a}}1 {ê}{{\^e}}1 {î}{{\^i}}1 {ô}{{\^o}}1 {û}{{\^u}}1%
    {Â}{{\^A}}1 {Ê}{{\^E}}1 {Ô}{{\^O}}1%
    {ç}{{\c{c}}}1 {Ç}{{\c{C}}}1%
    {ü}{{\"u}}1 {Ü}{{\"U}}1%
    {º}{{\textordmasculine}}1 {ª}{{\textordfeminine}}1%
}

% --- Diagramas (Cap. 2 hardware, Cap. 3 pipeline, Cap. 4 DFDs) ---
\usepackage{tikz}
\usetikzlibrary{
  shapes.geometric, arrows.meta, positioning, fit, calc, backgrounds,
  decorations.pathreplacing, shapes.multipart
}

% --- Equações/matemática leve ---
\usepackage{amsmath, amssymb}

% --- Compatibilização ABNT 6023:2018 e 10520:2023 ---
\usepackage{USPSC-classe/ABNT6023-10520}

% --- Links/hyperref ---
% A classe USPSC já configura hyperref; aqui apenas ajustamos as cores dos links.
\hypersetup{
  pdftitle={\TCCtitulo},
  pdfauthor={\TCCautor},
  pdfsubject={Monografia de Trabalho de Conclusão de Curso, Engenharia de Computação, ICMC-EESC/USP},
  pdfkeywords={visão computacional, monitoramento de fauna, gatos, USP São Carlos, edge AI, re-identificação},
  pdfcreator={LaTeX -- Pacote USPSC 3.2},
  colorlinks=true,
  linkcolor=blue!50!black,
  citecolor=blue!50!black,
  filecolor=magenta!50!black,
  urlcolor=blue!60!black,
  bookmarksdepth=4
}

% --- Caminhos para figuras ---
\graphicspath{{figuras/}{USPSC-img/}}

% --- Profundidade de sumário e numeração ---
% tocdepth=3 -> entra até subsubsection no sumário (omite \paragraph)
% secnumdepth=3 -> numera até subsubsection (omite numeração de \paragraph)
\setcounter{tocdepth}{3}
\setcounter{secnumdepth}{3}
